\bookchapter{נספח ג': טבלת השוואת “סלנג” מונחים מינימליסטית}

\begin{center}
\begin{tabularx}{\linewidth}{|X|X|X|}
\hline
\textbf{הסבר במשפט אחד} & \textbf{SoC / Chip} & \textbf{המוח והלב של הטלפון} \tabularnewline
\hline
הרכיב המרכזי שקובע אם הטלפון מהיר והמשחקים חלקים. & RAM (זיכרון) & "שולחן העבודה לטווח קצר" של הטלפון \tabularnewline
ככל שהשולחן גדול יותר, כך ניתן לפתוח יותר אפליקציות בו-זמנית והמעבר חלק יותר. & ROM (אחסון) & "המחסן הגדול" של הטלפון \tabularnewline
המקום שבו מאוחסנים תמונות, סרטונים ואפליקציות. 128G, 256G – זה מה שמתייחסים אליו. & OLED / LCD & שני סוגי חומרי מסך \tabularnewline
OLED – צבעים עשירים יותר, חסכוני יותר בחשמל; LCD – ידידותי יותר לעיניים, עלות נמוכה. & NFC & תקשורת אלחוטית לטווח קצר \tabularnewline
פונקציית ה"טאפ" שהופכת את הטלפון לכרטיס אוטובוס, כרטיס כניסה, כרטיס בנק. & OTA Update & עדכון מערכת אלחוטי \tabularnewline
ה"עדכון מערכת" הזה שלוחצים עליו ב[הגדרות], זה בדיוק זה. & Root & הרשאת מנהל-על של אנדרואיד \tabularnewline
Root & BL Lock & נעילת Bootloader \tabularnewline
מנעול אבטחה של היצרן, יש לבטלו לפני הבהוב. & NPU & מעבד רשת נוירונים \tabularnewline
שבב ייעודי למשימות AI, מהיר עשרות מונים מ-CPU, צריכת חשמל נמוכה יותר. & On-device AI & AI שרץ מקומית בטלפון \tabularnewline
AI רץ ישירות על הטלפון, לא צריך חיבור אינטרנט, מהיר, פרטיות טובה. & Qi2 & תקן טעינה אלחוטית מהדור החדש \tabularnewline
טעינה אלחוטית עם מיקום מגנטי, בדומה ל-MagSafe, טעינה נוחה יותר. & eSIM & כרטיס SIM מוטמע \tabularnewline
eSIM & LTPO & טכנולוגיית קצב רענון אדפטיבי \tabularnewline
LTPO & נספח ד': איפה למצוא משאבים אמינים? & תמיכה רשמית היא הבחירה הראשונה: \tabularnewline
\hline
\end{tabularx}
\end{center}
